\section*{Abstract}
	The Koide formula $Q = (m_e + m_\mu + m_\tau)/(\sqrt{m_e} + \sqrt{m_\mu} + \sqrt{m_\tau})^2 = 2/3$ is an empirical relation among charged lepton masses, confirmed experimentally to $0.001\%$, which the Standard Model cannot explain. In T0 theory, it emerges from the exponent structure $p = 3/2, 1, 2/3$ with step size $\Delta p = 1/3 = 1/D$ (three spatial dimensions of the torus manifold).
	
	We show that the \textbf{same} $1/3$-step generates the $g-2$ anomaly hierarchy $\Delta a_i = 4\pi/f^{q_i}$ with exponents $q = 5/3, 4/3$ --- again separated by $1/3$. The consequence is a direct, parameter-free bridge from the experimentally known mass ratio $m_\tau/m_\mu = 16.817$ to the $g$-$2$ ratio prediction $\Delta a(\tau{-}e)/\Delta a(\mu{-}e) = (144/125) \cdot m_\tau/m_\mu \approx 19.4$. This leads to $a_\tau \approx 1.28 \times 10^{-3}$ --- testable at Belle~II and qualitatively different from the SM prediction of $\approx 1.18 \times 10^{-3}$.

\section{Introduction: The Koide Mystery}

The charged lepton masses span three orders of magnitude:
\begin{equation}
	m_e = 0.511 \text{ MeV}, \quad m_\mu = 105.658 \text{ MeV}, \quad m_\tau = 1776.86 \text{ MeV}
\end{equation}

The Standard Model treats these as free parameters --- measured, not explained. Yet they satisfy the Koide relation~\cite{Koide1983} with extraordinary precision:
\begin{equation}
	\boxed{Q = \frac{m_e + m_\mu + m_\tau}{\left(\sqrt{m_e} + \sqrt{m_\mu} + \sqrt{m_\tau}\right)^2} = 0.666661 \approx \frac{2}{3}}
	\label{eq:koide}
\end{equation}
confirmed to $\Delta Q/Q = 0.001\%$. This precision demands explanation, yet the SM offers none. It is as if the masses know about each other --- they satisfy a collective relation that no accident could produce.

\section{T0 Explanation of Koide: The Exponent Structure}

In T0 theory, all lepton masses derive from a single parameter $\xi = 4/30{,}000$ via the Yukawa structure~\cite{Pascher006}:
\begin{equation}
	m_i = r_i \cdot \xi^{p_i} \cdot v, \quad v = 246 \text{ GeV (Higgs VEV)}
	\label{eq:yukawa}
\end{equation}
with rational prefactors $r_i$ and exponents $p_i$:

\begin{table}[h]
	\centering
	\begin{tabular}{lccccc}
		\toprule
		\textbf{Lepton} & $r_i$ & $p_i$ & $m_{\text{T0}}$ (MeV) & $m_{\text{exp}}$ (MeV) & $\Delta\%$ \\
		\midrule
		$e$   & $4/3$  & $3/2$ & $0.505$  & $0.511$  & $-1.2\%$ \\
		$\mu$ & $16/5$ & $1$   & $104.96$ & $105.66$ & $-0.7\%$ \\
		$\tau$& $25/9$ & $2/3$ & $1783.4$ & $1776.9$ & $+0.4\%$ \\
		\bottomrule
	\end{tabular}
	\caption{T0 lepton masses from a single parameter $\xi$}
	\label{tab:masses}
\end{table}

The exponents form an arithmetic sequence:
\begin{equation}
	\boxed{p_e = \frac{3}{2}, \quad p_\mu = 1, \quad p_\tau = \frac{2}{3}, \quad \Delta p = \frac{1}{3}}
	\label{eq:exponents}
\end{equation}

This uniform step of $1/3$ is not arbitrary. In T0 theory, the universe has the topology of a four-dimensional torus with three spatial dimensions ($D = 3$). The lepton generations correspond to different winding modes on this torus, and the step between generations is $1/D = 1/3$.

The Koide relation then follows naturally: the specific combination of rational prefactors $(4/3, 16/5, 25/9)$ and the geometric progression $\xi^{3/2}, \xi^1, \xi^{2/3}$ conspire to produce $Q \approx 2/3$. What Koide discovered empirically, T0 derives from the torus geometry~\cite{Pascher116}.

\section{From Masses to $g-2$: The Same Step Size}

\subsection{The $g-2$ Formulas in T0}

The anomalous magnetic moments in T0 theory are~\cite{Pascher018}:
\begin{align}
	a_e &= \frac{S_3}{f \cdot k_\text{eff}}, \quad S_3 = 2\pi^2, \quad f = 1/\xi = 7500 \\
	a_\mu &= a_e + \frac{4\pi}{f^{5/3}} \\
	a_\tau &= a_e + \frac{4\pi}{f^{4/3}}
\end{align}

The lepton-specific corrections $\Delta a_i$ above the electron baseline have the structure:
\begin{equation}
	\Delta a_i = \frac{4\pi}{f^{q_i}}, \quad q_\mu = \frac{5}{3}, \quad q_\tau = \frac{4}{3}
	\label{eq:g2exponents}
\end{equation}

\subsection{The Key Observation: Same $1/3$ Step}

The $g-2$ exponents:
\begin{equation}
	q_\mu = \frac{5}{3}, \quad q_\tau = \frac{4}{3}, \quad \Delta q = \frac{1}{3}
\end{equation}

The mass exponents (Eq.~\ref{eq:exponents}):
\begin{equation}
	p_\mu = 1, \quad p_\tau = \frac{2}{3}, \quad \Delta p = \frac{1}{3}
\end{equation}

\begin{equation}
	\boxed{\Delta p = \Delta q = \frac{1}{3} = \frac{1}{D}}
\end{equation}

This is the central result: \textbf{The step between lepton generations is $1/3$ in both the mass hierarchy and the $g-2$ hierarchy.} This is not a fit --- it is the same geometric origin (three spatial dimensions of the torus) manifesting in two different physical observables.

\subsection{Why This Must Be So}

In T0 theory, both mass and anomalous moment arise from how a lepton couples to the torus winding structure:
\begin{itemize}
	\item The \textbf{mass} measures the coupling strength to the vacuum ($\xi$-field), governed by $\xi^{p_i}$
	\item The \textbf{anomalous moment} measures the deviation from the Dirac value due to the geometric environment, governed by $f^{-q_i}$
\end{itemize}
Both are aspects of the same winding mode. The step $1/3$ between generations reflects the dimensionality of the spatial manifold. If the torus had $D = 4$ spatial dimensions, the step would be $1/4$.

\section{The Bridge: Mass Ratios $\to$ $g-2$ Ratios}

\subsection{The $f^{1/3}$ Factor}

From the mass formulas:
\begin{equation}
	\frac{m_\tau}{m_\mu} = \frac{r_\tau}{r_\mu} \cdot \xi^{p_\tau - p_\mu} 
	= \frac{25/9}{16/5} \cdot \xi^{-1/3} = \frac{125}{144} \cdot f^{1/3}
	\label{eq:massratio}
\end{equation}

From the $g-2$ formulas:
\begin{equation}
	\frac{\Delta a(\tau{-}e)}{\Delta a(\mu{-}e)} 
	= \frac{4\pi/f^{4/3}}{4\pi/f^{5/3}} = f^{1/3}
	\label{eq:g2ratio}
\end{equation}

The factor $f^{1/3}$ appears in both. Combining Eqs.~\ref{eq:massratio} and~\ref{eq:g2ratio}:
\begin{equation}
	\boxed{\frac{\Delta a(\tau{-}e)}{\Delta a(\mu{-}e)} = \frac{144}{125} \cdot \frac{m_\tau}{m_\mu}}
	\label{eq:bridge}
\end{equation}

This is the bridge: \textbf{The $g-2$ ratio is determined by the mass ratio}, up to the known rational factor $144/125 = 1.152$. No free parameters. No $\alpha$. No $k_\text{eff}$. No $K_\text{frak}$. Everything cancels except the mass ratio and the rational prefactors.

\subsection{Numerical Verification}

Using the experimentally known mass ratio:
\begin{align}
	\frac{m_\tau}{m_\mu} &= \frac{1776.86}{105.658} = 16.817 \\
	\frac{144}{125} \times 16.817 &= 19.373
\end{align}

Directly from $f^{1/3}$:
\begin{equation}
	f^{1/3} = 7500^{1/3} = 19.574
\end{equation}

The $1.0\%$ difference between 19.373 and 19.574 reflects the small deviation of T0 masses from experiment (Table~\ref{tab:masses}). Both values are far from the SM ratio (see Section~\ref{sec:comparison}).

\subsection{Parameter-Free Tau Prediction}

Applying Eq.~\ref{eq:bridge} with experimental values only:
\begin{align}
	\Delta a(\mu{-}e)_\text{exp} &= a_\mu^\text{exp} - a_e^\text{exp} = 6.27 \times 10^{-6} \\
	\Delta a(\tau{-}e) &= \Delta a(\mu{-}e)_\text{exp} \times \frac{144}{125} \times \frac{m_\tau}{m_\mu} \\
	&= 6.27 \times 10^{-6} \times 19.37 = 1.214 \times 10^{-4} \\
	a_\tau &= a_e^\text{exp} + \Delta a(\tau{-}e) \\
	&= 1.160 \times 10^{-3} + 1.214 \times 10^{-4} \\
	&= \boxed{1.281 \times 10^{-3}}
	\label{eq:taupred}
\end{align}

This prediction uses:
\begin{itemize}
	\item Experimental $a_e$ and $a_\mu$ (no T0 absolute values)
	\item Experimental $m_\tau/m_\mu$ (no T0 mass formulas)
	\item Only the rational factor $144/125$ from the T0 prefactors
\end{itemize}
No $\alpha$, no $K_\text{frak}$, no $k_\text{eff}$. Maximum predictive purity.

\section{Comparison with the Standard Model}
\label{sec:comparison}

\subsection{Mass Ratios: SM Has No Prediction}

The SM treats all lepton masses as free parameters. It cannot predict $m_\tau/m_\mu$, cannot explain Koide, and has no mass hierarchy formula. The experimental ratios $m_\mu/m_e = 206.77$, $m_\tau/m_\mu = 16.82$ are inputs, not outputs.

T0 derives all three masses from $\xi$, explains Koide, and predicts the ratios to $\sim 1\%$.

\subsection{$g-2$ Ratios: Qualitatively Different Predictions}

The SM calculates each lepton's $g-2$ separately using perturbative Feynman diagrams (up to 12,672 at fifth order for the electron). No scaling relation between leptons is built in --- the ratio of anomalies simply falls out of the calculation.

\begin{table}[h]
	\centering
	\begin{tabular}{lccc}
		\toprule
		\textbf{Ratio} & \textbf{T0} & \textbf{SM} & \textbf{Factor} \\
		\midrule
		$\Delta a(\tau{-}e) / \Delta a(\mu{-}e)$ & $19.57$ & $2.80$ & $7.0\times$ \\
		$a_\tau$ prediction & $1.28 \times 10^{-3}$ & $1.18 \times 10^{-3}$ & $+8.5\%$ \\
		\bottomrule
	\end{tabular}
	\caption{T0 vs SM: $g-2$ ratio predictions for the tau lepton}
	\label{tab:g2comparison}
\end{table}

This is not a subtle $0.1\%$ difference. The T0 $g-2$ ratio is \textbf{seven times larger} than the SM ratio. The tau anomaly differs by $8.5\%$. Belle~II will decide.

\subsection{The Fair Comparison: Equal Starting Conditions}

A common objection is that SM calculations of $a_e$ achieve $10^{-12}$ precision, far beyond T0. This comparison is however misleading:

\begin{itemize}
	\item The SM extracts $\alpha$ from the $a_e$ measurement itself. Calculating $a_e$ with this $\alpha$ is circular --- it is a fit, not a prediction.
	\item T0 derives $\alpha \approx 4\pi^3/(f \cdot k_\text{eff})$, yielding $1/\alpha \approx 137.22$ ($0.14\%$ from experiment) with no input.
	\item When the SM uses $\alpha_\text{T0}$ instead of $\alpha_\text{exp}$, its $a_\mu$ prediction deviates by $-0.137\%$ --- comparable to T0's own $-0.148\%$ deviation. At equal footing, the 12,672 diagrams offer no advantage over T0's single geometric formula.
\end{itemize}

The ratio prediction (Eq.~\ref{eq:bridge}) sidesteps this debate entirely: it uses experimental $a_e$, $a_\mu$, and $m_\tau/m_\mu$. The only theoretical input is the rational factor $144/125$.

\section{The Argument Chain}

The logic is:
\begin{enumerate}
	\item The Koide formula holds experimentally to $0.001\%$. This is an unexplained empirical fact.
	\item T0 explains Koide through the exponent structure $p = 3/2, 1, 2/3$ with step $1/3 = 1/D$.
	\item The \textbf{same} $1/3$ step appears in the $g-2$ exponents $q = 5/3, 4/3$.
	\item This is not a coincidence --- it is the same torus winding structure manifesting in two observables.
	\item Therefore: the $g-2$ ratio $\Delta a(\tau{-}e)/\Delta a(\mu{-}e) = f^{1/3}$ is as well-founded as the mass hierarchy. It rests on the same geometry.
	\item Combining with experimental mass ratios: $a_\tau \approx 1.28 \times 10^{-3}$.
\end{enumerate}

If one accepts that the Koide relation is not accidental --- and its $0.001\%$ precision makes accident implausible --- then one must also accept that the $g-2$ ratio prediction has the same geometric foundation. The predictions are linked. They cannot be separated.

\section{Belle~II: The Decisive Test}

Belle~II at KEK expects sensitivity of $\sim 10^{-7}$ for $a_\tau$. The predictions diverge by $8.5\%$:

\begin{table}[h]
	\centering
	\begin{tabular}{lc}
		\toprule
		\textbf{Source} & $a_\tau$ ($\times 10^{-3}$) \\
		\midrule
		SM (perturbative, $\alpha_\text{exp}$ input) & $1.177$ \\
		T0 geometric (corrected) & $1.245$ \\
		T0 from mass ratio (Eq.~\ref{eq:bridge}) & $1.281$ \\
		T0 from $f^{1/3}$ (Eq.~\ref{eq:g2ratio}) & $1.282$ \\
		\midrule
		T0 range & $1.25$--$1.28$ \\
		\bottomrule
	\end{tabular}
	\caption{Tau anomaly predictions --- Belle~II will distinguish}
\end{table}

Three outcomes are possible:
\begin{itemize}
	\item $a_\tau \approx 1.18 \times 10^{-3}$: SM confirmed, T0 falsified.
	\item $a_\tau \approx 1.25$--$1.28 \times 10^{-3}$: T0 confirmed. The mass-$g$-$2$ bridge is real, and the Koide structure extends to anomalous moments.
	\item $a_\tau$ outside both ranges: both theories require revision.
\end{itemize}

\section{Conclusion}

The Koide formula is not an isolated curiosity. In T0 theory, it is one manifestation of a deeper geometric structure: the $1/3$-step between lepton generations, originating from the three-dimensional torus topology. This same structure predicts the $g-2$ anomaly hierarchy, connecting two seemingly unrelated observables through a single geometric principle.

The mass ratio $m_\tau/m_\mu = 16.82$ --- experimentally known to high precision --- directly determines the $g-2$ ratio $\Delta a(\tau{-}e)/\Delta a(\mu{-}e) \approx 19.4$, leading to $a_\tau \approx 1.28 \times 10^{-3}$. This prediction is parameter-free, $\alpha$-independent, and qualitatively different from the SM value of $1.18 \times 10^{-3}$.

The Standard Model cannot make this connection. It has no explanation for Koide, no scaling relation between generations, and no bridge from masses to anomalous moments. It treats each observable as independent. T0 treats them as two faces of one geometry.

Belle~II will decide.

\begin{thebibliography}{9}
	\bibitem{Koide1983} Y.~Koide, ``New viewpoint on quark and lepton mass matrix'', 
	\textit{Phys.\ Rev.\ Lett.}\ \textbf{47}, 1241 (1983).
	
	\bibitem{Pascher006} J.~Pascher, ``Yukawa Couplings and Mass Generation in T0 Theory'', 
	Document 006, T0-Time-Mass-Duality Repository (2025).
	\url{https://github.com/jpascher/T0-Time-Mass-Duality/blob/main/2/pdf/}
	
	\bibitem{Pascher116} J.~Pascher, ``The Koide Formula as a Consequence of T0 Theory'', 
	Document 116, T0-Time-Mass-Duality Repository (2025).
	\url{https://github.com/jpascher/T0-Time-Mass-Duality/blob/main/2/pdf/}
	
	\bibitem{Pascher018} J.~Pascher, ``Anomalous Magnetic Moments in T0 Theory --- Complete Geometric Derivation with Fractal Correction'', 
	Document 018 Rev.~11, T0-Time-Mass-Duality Repository (2026).
	\url{https://github.com/jpascher/T0-Time-Mass-Duality/blob/main/2/pdf/}
	
	\bibitem{EidelmanPassera2007} S.~Eidelman, M.~Passera, ``Tau anomalous magnetic moment'', 
	\textit{Mod.\ Phys.\ Lett.\ A} \textbf{22}, 159 (2007).
\end{thebibliography}
